% !PRINTED 490
\weberpage{490}

An algebraic number satisfies not merely one, but infinitely many rational equations; for if we multiply two arbitrary functions of the form $F(x)$ with one another, we obtain a function of the same form, which vanishes for $x = \Theta$ if one of the factors has this property.

Among all the rational equations that an algebraic number satisfies there is one of lowest possible degree, $f(\Theta) = 0$, where $f(x)$ has the form:
\weq{2}{f(x) = x^n + a_1 x^{n-1} + a_2 x^{n-2} + \cdots + a_n,}
and there can also be only one such equation, if, as always hitherto, we take the coefficient of the highest power of $x$ equal to 1.

For if $f(x), f_1(x)$ are two functions of the form (2) of the same degree $n$, then $f(x) - f_1(x)$ is of degree lower than $n$, and if both $f(\Theta)$ and $f_1(\Theta)$ vanish, then $f(\Theta) - f_1(\Theta)$ vanishes also; if, then, this difference does not vanish identically, $\Theta$ satisfies an equation of degree lower than $n$, which is contrary to the assumption.

\gesperrt{The function $f(x)$ is irreducible in the field of rational numbers.}

For if $f(x)$ splits into two rational factors $f_1(x)$ and $f_2(x)$, each of which is of lower degree than $f(x)$, then $\Theta$ satisfies one of the two equations $f_1(\Theta) = 0$, $f_2(\Theta) = 0$, which contradicts our assumption.

\gesperrt{If $n$ is the degree of the rational equation of lowest degree that the number $\Theta$ satisfies, we call $\Theta$ an algebraic number of degree $n$.}

\begin{center}
\S.\,133.

\gesperrt{Integral algebraic numbers.}
\end{center}

\gesperrt{An algebraic number $\Theta$ is called an integral algebraic number if it satisfies a rational equation}
\weq{1}{\Theta^m + A_1 \Theta^{m-1} + \cdots + A_{m-1} \Theta + A_m = 0}
%CONT>>
